% Springer Lecture Notes in Computer Science (LNCS) proceedings.
\documentclass[runningheads]{llncs}

\usepackage[T1]{fontenc}
\usepackage{graphicx}
\usepackage{amsmath}

\begin{document}

\title{An Example Title}
\titlerunning{An Example Title}

\author{A. U. Thor\inst{1}\orcidID{0000-0000-0000-0000} \and
Second Author\inst{2}}
\authorrunning{A. U. Thor and S. Author}

\institute{Your University, City, Country\\
\email{author@example.org} \and
Another Institute, City, Country}

\maketitle

\begin{abstract}
State the problem, the approach and the main result in 150 to 250 words.

\keywords{First keyword \and Second keyword \and Third keyword.}
\end{abstract}

\section{Introduction}

Say what is known, what is missing, and what this paper adds. Cite earlier
work by number \cite{einstein1905,knuth1984}.

\section{Method}

Describe what you did in enough detail to repeat it.

\begin{equation}
  \bar{x} = \frac{1}{n} \sum_{i=1}^{n} x_i
  \label{eq:mean}
\end{equation}

\begin{table}
  \caption{Example measurements.}
  \label{tab:results}
  \centering
  \begin{tabular}{|l|c|c|}
    \hline
    Sample & Mean & Std.\ dev. \\
    \hline
    A & 1.02 & 0.11 \\
    B & 0.97 & 0.08 \\
    \hline
  \end{tabular}
\end{table}

\section{Results}

Report what you found. Table~\ref{tab:results} lists the measurements and
Eq.~(\ref{eq:mean}) defines the mean.

\begin{credits}
\subsubsection{\ackname} Thank the people and grants that made the work possible.
\end{credits}

\bibliographystyle{splncs04}
\bibliography{references}

\end{document}
